\chapter{基于深度射频指纹的复杂环境单目标高精度定位方法}
\label{ch:single_target}

\section{问题描述与总体思路}

城市峡谷环境下的无线传播呈现出极强的非平稳特性。高层建筑的密集分布与街道空间的狭窄几何结构，使得低轨卫星下行信号在到达接收端前，普遍经历了复杂的反射、绕射及散射过程。这种典型的非视距（NLOS）传播导致接收端观测到的信号由多条传播路径叠加而成，破坏了传统几何定位方法（如基于 RSSI 或 AoA）赖以生存的直达径假设。

在复杂多径场景下，RSSI 易受阴影衰落与干涉波动的剧烈影响，AoA 测量值常因强反射分量而偏离真实方位，而 TOA/TDOA 则会因为首径被遮挡产生显著的距离正偏差。这种误差并非来自单纯的随机噪声，而是传播环境直接改变了物理观测量的语义信息。

然而，多径效应在干扰几何测量的同时，也赋予了信号丰富的位置相关性。不同空间坐标下的建筑布局与反射面组合，在信号时域和频域上留下了独特的“传播签名”。本章核心思路是将这些复杂的多径分量视为可利用的特征，而非干扰。通过构建高保真的射线追踪仿真环境，重构包含多径衰落特征的 I/Q 双路复基带序列，并利用深度残差网络（ResNet-50）从中自动提取高维射频指纹特征。相比人工构造的低维观测量，这种深度射频指纹能更完整地捕捉信号在幅度、相位及多径时延结构上的微小差异，从而实现复杂场景下的高精度定位。

\begin{figure}[htbp]
    \centering
    % \includegraphics[width=0.85\textwidth]{figures/ch3_workflow.pdf}
    \vspace{4cm}
    \caption{基于深度射频指纹的单目标定位方法总体架构}
    \label{fig:ch3_overall_architecture}
\end{figure}

\section{低轨卫星信号特征构造与样本生成}

\subsection{Starlink 下行链路帧结构分析}

为使信号建模过程与低轨卫星通信场景保持一致，本文选取 Starlink Ku 波段下行链路作为物理层建模对象。建模的重点在于提取与定位特征相关的时频结构参数，而非复现完整的通信协议。

Starlink 下行链路采用 OFDM 体制。根据公开文献分析，单信道带宽通常为 $240\,\text{MHz}$，子载波数为 $1024$。单个 OFDM 符号包含 $4.266\,\mu\text{s}$ 的有用时长及 $0.133\,\mu\text{s}$ 的保护间隔时长。帧周期约为 $1/750\,\text{s}$，每帧包含 302 个非零符号。此外，低轨卫星的高速运动带来的多普勒（Doppler）效应对信号时域对齐及 I/Q 采样值的相位旋转具有显著影响。本章在样本生成阶段，将卫星相对运动产生的频移与时间压缩效应显式引入传播模型，以保证射频指纹样本的物理一致性。

为避免信号建模参数在正文中分散出现，本文将参考信号构造过程中采用的主要参数统一列于表~\ref{tab:signal_params_ch3}。

\begin{table}[htbp]
    \centering
    \bicaption{低轨卫星下行参考信号建模参数设置}{Modeling parameters for LEO satellite downlink reference signal}
    \label{tab:signal_params_ch3}
    \begin{tabular}{@{}llccl@{}}
        \toprule
        \textbf{参数类别} & \textbf{参数名称} & \textbf{符号} & \textbf{取值/设置} & \textbf{说明} \\ \midrule
        频域参数 & 载波中心频率 & $f_c$ & Ku 波段相关频率 & 参考信号载频设置 \\
        & 单信道带宽 & $F_s$ & $240\,\text{MHz}$ & 单信道有效带宽 \\
        & 子载波数 & $N$ & 1024 & OFDM 参考参数 \\
        & 循环前缀长度 & $N_g$ & 32 & OFDM 参考参数 \\
        时域参数 & 有效符号时长 & $T$ & $4.266\,\mu\text{s}$ & $T=N/F_s$ \\
        & 保护间隔时长 & $T_g$ & $0.133\,\mu\text{s}$ & $T_g=N_g/F_s$ \\
        & 总符号时长 & $T_{\mathrm{sym}}$ & $4.4\,\mu\text{s}$ & $T_{\mathrm{sym}}=T+T_g$ \\
        帧参数 & 帧周期 & $T_f$ & $1/750\,\text{s}$ & 参考帧级时间结构 \\
        & 每帧非零符号数 & — & 302 & 帧内数据符号数 \\
        频谱结构 & 信道数量 & — & 8 & 多信道组织方式 \\
        & 相邻信道间隔 & — & $250\,\text{MHz}$ & 频谱布局参数 \\
        & 保护带宽 & — & $10\,\text{MHz}$ & 信道间保护带 \\
        & 中心 gutter 宽度 & — & 4 个子载波 & 减小中心杂散影响 \\
        同步结构 & 主同步字段 & PSS & 保留 & 同步锚定与特征构造 \\
        & 辅同步字段 & SSS & 保留 & 同步与结构对齐 \\
        & 其他辅助字段 & CSS/CM1SS & 按需选取 & 仅保留定位相关结构 \\
        传播参数 & 多普勒频移 & $f_d$ & 场景相关 & 星地相对运动决定 \\
        信道参数 & 有效传播路径数 & $L$ & 场景相关 & 射线追踪仿真结果 \\
        噪声参数 & 信噪比范围 & SNR & $[-15, 20]\,\text{dB}$ & 鲁棒性实验设置 \\
        样本形式 & 输入样本表示 & $\mathbf{X}_{\mathrm{in}}$ & I/Q 双通道矩阵 & 用于射频指纹学习 \\ \bottomrule
    \end{tabular}
\end{table}

\subsection{I/Q 双路复基带信号重构}

本文直接采用完整 I/Q 双路复基带采样作为定位模型输入。设发射端参考基带信号为 $s(t)$，在包含 $L$ 条有效传播路径的信道下，接收端的复基带信号 $r(t)$ 可建模为：
\begin{equation}
r(t) = \sum_{l=1}^{L} \beta_l \cdot G(\theta_l, \varphi_l) \cdot s(t - \tau_l) e^{j2\pi f_{d,l} t} + n(t)
\end{equation}
其中，$\beta_l$ 为第 $l$ 条路径的复增益，$G(\theta_l, \varphi_l)$ 为接收天线增益加权。

在数字下变频后，将其拆分为实部（I路）与虚部（Q路），构造出维度为 $2 \times N$ 的输入矩阵：
\begin{equation}
\mathbf{X}_{\mathrm{in}} = \begin{bmatrix} 
\Re\{r[1]\} & \Re\{r[2]\} & \dots & \Re\{r[N]\} \\ 
\Im\{r[1]\} & \Im\{r[2]\} & \dots & \Im\{r[N]\} 
\end{bmatrix}
\end{equation}

\subsection{样本标注与预处理}

本文利用三维仿真在目标区域内生成了覆盖 $27,546$ 个网格位置的样本集。位置标签采用平面坐标 $(x_k, y_k)$。对输入张量执行基于通道维度的标准化处理：
\begin{equation}
\tilde{\mathbf{X}}_{\mathrm{in}}^{i,j} = \frac{\mathbf{X}_{\mathrm{in}}^{i,j} - \mu_i}{\sigma_i}
\end{equation}

\section{基于 ResNet-50 的射频指纹模型设计}

\subsection{网络架构设计}

本章采用具有深层表征能力的 ResNet-50 作为主干网络。其核心残差结构定义为：
\begin{equation}
\mathbf{y} = \mathcal{F}(\mathbf{x}) + \mathbf{x}
\end{equation}
通过跳跃连接，模型能够稳定地从高采样率信号中学习深层的物理指纹。

\subsection{损失函数与训练策略}

本文采用多任务学习架构。总损失函数定义为：
\begin{equation}
\mathcal{L}_{\mathrm{total}} = \lambda_1 \mathcal{L}_{\mathrm{cls}} + \lambda_2 \mathcal{L}_{\mathrm{reg}}
\end{equation}

本文采用的 ResNet-50 网络结构及训练配置如表~\ref{tab:resnet_params_ch3} 所示。该表同时给出了输入形式、主干网络配置、损失函数及训练参数设置。

\begin{table}[htbp]
    \centering
    \bicaption{ResNet-50 射频指纹定位模型结构与训练参数设置}{Model structure and training parameters of ResNet-50 for RF fingerprinting}
    \label{tab:resnet_params_ch3}
    \begin{tabular}{@{}lll@{}}
        \toprule
        \textbf{参数类别} & \textbf{参数名称} & \textbf{取值/设置} \\ \midrule
        输入参数 & 输入形式 & I/Q 双通道矩阵 \\
        & 输入尺寸 & $2 \times N$ ($N$ 为时域采样点数) \\
        主干网络 & 主干网络 (Backbone) & ResNet-50 \\
        & 残差阶段配置 & $\{3, 4, 6, 3\}$ \\
        & 池化方式 & 全局平均池化 (Global Average Pooling) \\
        & 输出特征维度 & 2048 \\
        分支任务 & 分类分支输出数 & $C$ (由区域划分数决定) \\
        & 回归分支输出维度 & 2 (对应平面坐标 $(x, y)$) \\
        损失函数 & 分类损失函数 & 交叉熵损失 (Cross-Entropy) \\
        & 回归损失函数 & 均方误差 (MSE) \\
        & 联合损失形式 & $\lambda_1 \mathcal{L}_{\mathrm{cls}} + \lambda_2 \mathcal{L}_{\mathrm{reg}}$ \\
        优化策略 & 优化器 (Optimizer) & Adam \\
        & 批大小 (Batch Size) & 128 \\
        & 训练轮次 (Epoch) & 150 \\
        & 学习率策略 (LR Schedule) & 衰减策略 \\
        数据集划分 & 划分比例 (Train:Val:Test) & 8:1:1 \\
        评价指标 & 性能评价指标 & F1-score, RMSE, Inference Time \\
        扩展支持 & 多天线输入支持 & 支持 (可扩展至 MIMO 多通道) \\ \bottomrule
    \end{tabular}
\end{table}

\section{MIMO 空间分集增强机制}

为突破低信噪比环境下的定位瓶颈，本章引入 MIMO 空间分集技术。通过在接收端部署 $K$ 通道阵列，采集 $K \times N$ 的时空矩阵：
\begin{equation}
\mathbf{R}_{\mathrm{MIMO}}[n] = \big[ r_1[n], r_2[n], \dots, r_K[n] \big]^T
\end{equation}
多通道输入使得网络能够在隐空间中自发实现能量聚合，有效对抗信道深度衰落。

\section{实验验证与性能分析}

\subsection{定位精度对比分析}

为比较不同方法在复杂城市环境下的单目标定位性能，本文将几何特征方法、浅层卷积网络和深层残差网络的测试结果统一列于表~\ref{tab:perf_comparison_ch3}。

\begin{table}[htbp]
    \centering
    \bicaption{不同单目标定位方法性能对比}{Performance comparison of different single-target localization methods}
    \label{tab:perf_comparison_ch3}
    \begin{tabular}{@{}llcccc@{}}
        \toprule
        \textbf{方法} & \textbf{输入特征} & \textbf{模型类型} & \textbf{RMSE / m} & \textbf{F1-score} & \textbf{推理时间 / ms} \\ \midrule
        Geometric-based & RSSI, AoA 等 & 传统回归模型 & $45.2$ & -- & -- \\
        4-Layer CNN & I/Q 样本 & 浅层卷积网络 & 待补 & $0.73$ & -- \\
        ResNet-50 & I/Q 样本 & 深层残差网络 & $\mathbf{8.3}$ & $\mathbf{0.85}$ & $\mathbf{85.0}$ \\
        ResNet-101 & I/Q 样本 & 更深层残差网络 & 待补 & 待补 & $142.0$ \\ \bottomrule
    \end{tabular}
\end{table}

\subsection{计算开销与鲁棒性折衷}

分析表明 ResNet-50 在保持高精度的同时，单帧推理时延仅为 $85.0\,\text{ms}$，满足实时性需求。而 ResNet-101 推理时延过长，不利于工程部署。

\subsection{极端信噪比条件下的鲁棒性评估}

针对城市建筑物对星-地链路的高额传输损耗，本节测试了系统在 $\text{SNR} = -15\,\text{dB}$ 环境下的表现。

\subsection{不同 MIMO 配置下的性能分析}

为进一步分析空间分集机制对低信噪比场景定位性能的影响，本文将不同接收阵列配置下的实验结果列于表~\ref{tab:mimo_comparison_ch3}。

\begin{table}[htbp]
    \centering
    \bicaption{不同 MIMO 配置下低信噪比场景定位性能对比}{Localization performance with different MIMO configurations in low SNR scenarios}
    \label{tab:mimo_comparison_ch3}
    \begin{tabular}{@{}lcccl@{}}
        \toprule
        \textbf{接收配置} & \textbf{输入通道数} & \textbf{典型 SNR 条件} & \textbf{RMSE / m} & \textbf{性能描述} \\ \midrule
        SISO & 2 & $-15\,\text{dB}$ & 待补 & 误差显著增大 \\
        $2\times 2$ MIMO & 4 & $-15\,\text{dB}$ & 待补 & 较 SISO 有明显改善 \\
        $4\times 4$ MIMO & 8 & $-15\,\text{dB}$ & 待补 & 分集进一步降低误差 \\
        $8\times 8$ MIMO & 16 & $-15\,\text{dB}$ & $\mathbf{< 20}$ & 保持较好鲁棒性 \\ \bottomrule
    \end{tabular}
\end{table}

\section{本章小结}

本章通过融合深度射频指纹特征提取与 MIMO 空间分集增强技术，显著提升了复杂城市环境下单目标定位的精度与鲁棒性。实验证明，所提方法在保持实时性的前提下，即使在极端低信噪比环境下也能提供可靠的定位结果。
